\section{共享年代之二：新朝、复汉与名义的分裂}

新朝改钱、改官名、改革土地和爵制，留下钱币、官印和《汉书·王莽传》可供交叉的痕迹；它们不能证明每道命令以同一速度进入乡里和边郡。制度名词越频繁变化，越能提示书面命名与实际组织之间存在距离。天灾、价格、地方冲突与战争如何彼此作用，材料不足以编成一个单因果故事；较稳妥的是说，中枢的命令链正在失去使各处相信、供给和执行的条件。

昆阳、绿林、赤眉与更始政权的细节各有史书立场，不能被一条“农民起义”的标签抹平。二五年，\eventref{EH-0025-GUANGWU-ACCESSION}{刘秀在鄗称帝}，又\eventref{EH-0025-LUOYANG-CAPITAL}{定都雒阳}。这个选择靠近关东、河北的道路和联盟，也承认长安周边已难稳定承担大军供给。二六年赤眉入长安，更始政权终结；二七年\eventref{EH-0027-CHIMEI-DEFEAT}{赤眉主力受降}。陇右、巴蜀与岭南的军事中心随后才逐步被接回雒阳的诏令网络，“平定”从来不只是地图重新着色。

建武度田说明复兴如何进入地方日常：朝廷试图重写战后土地与人户，地方反弹和隐匿争议说明簿籍不是透明事实。交趾的\eventref{EH-0040-TRUNG-REBELLION}{徵氏起兵}与前往南方的\eventref{EH-0043-MA-YUAN-JIAOZHI}{马援军}，也不能只讲作一场边地被“完全纳入”的战争。汉文正史与后来的越南叙事对人物及意义有不同安排；它们至少共同提示，河网、任命与地方首领关系是长久的行政问题。

北边和西域让复兴的半径更显有限。七三年\eventref{EH-0073-DOU-GU-YIWULU}{伊吾卢与班超赴鄯善}，九一年\eventref{EH-0091-BAN-CHAO-PROTECTOR}{班超为都护}，是使节、驻军、绿洲王权与河西供给重新接合的节点，而不是稳定疆界的证明。九七年甘英到达安息属地，\eventref{EH-0097-GAN-YING-ANXI}{未至大秦}；所闻的海路与远方信息可说明知识传播，不能改写为汉罗直接外交。

九十年代的宫廷转折把军功和宫门连在一起。八九年窦宪北出并立燕然山铭，\eventref{EH-0089-DOU-XIAN-JILU}{纪功文字}须和持久控制、兵数分开；九二年\eventref{EH-0092-DOU-CLAN-PURGE}{和帝与郑众诛窦氏}，解决了外戚的即时军政主导，也使宦官更靠近传诏与宫廷安全。此后外戚、宦官、幼帝和官僚的反复冲突，不能降格为人格故事：他们争夺的是禁军、诏令、人事和信息的高风险节点。

一六六年和一六九年的\eventref{EH-0166-FIRST-PARTISAN-PROHIBITION}{两次党锢}，表明朝廷将部分太学生、官僚和地方声望网络视为竞争资源；名单、地域和结社程度有争议，不能把“党人”写成一致的组织。一七五年的\eventref{EH-0175-XIPING-STONE-CLASSICS}{熹平石经}又显示朝廷仍在太学公开校定可援引的经典。压制与定经并置，正好显示制度同时需要公共评价，又害怕评价形成独立渠道。

一八四年\eventref{EH-0184-YELLOW-TURBAN}{黄巾起事}在多州郡展开。饥馑、赋役、地方保护和宗教语言在不同地区构成不同条件；主力被击破不等于乡里压力消失。朝廷为多战线征调，地方州郡、边将与豪强因而获得更多招募、供给和谈判空间。后来割据不能倒写成一八八年设置州牧时唯一且确定的意图，但州的军政用途确实扩展了。

一八九年灵帝死后，\eventref{EH-0189-HE-JIN-MASSACRE}{何进被杀、宦官遭屠戮}，少帝和陈留王一度离京；责任先后、董卓是否被谁召入和伤亡人数都有异说。董卓随后\eventref{EH-0189-DONG-ZHUO-DEPOSITION}{废少帝立献帝}，次年\eventref{EH-0190-CHANGAN-TRANSFER}{西迁长安}。此后问题不再是“谁有皇帝”，而是谁能让皇帝的诏令经过粮仓、道路、军队和地方官而成为行动。

曹操一九六年\eventref{EH-0196-CAO-CAO-RECEIVES-EMPEROR}{迎献帝至许}，以屯田、官吏和军队把名义接入北方资源；屯田规模和覆盖地区仍须逐案考察。二〇〇年\eventref{EH-0200-GUANDU}{官渡胜利}依赖乌巢粮仓、情报、部属选择及河北、河南的供给条件，不是单一计策的结果。二〇八年\eventref{EH-0208-RED-CLIFFS}{赤壁受挫}则显示长江、水军和地方联盟限制了北方力量的南下。

二二〇年曹操去世，曹丕接管魏王位和朝廷军政；十月\eventref{EH-0220-HAN-ABDICATION}{献帝禅位}，东汉皇帝继承体系结束。诏册能证明仪式被完成，不能证明权力转移出于自愿，也不能解释蜀、吴为何不接受魏的统一权威。汉的接口没有消失：它们被魏、蜀、吴以不同方式继承、争夺和改造。

\begin{sourcenote}[复汉与禅代的两个陷阱]
《后汉书》常从东汉正统回望王莽与复汉，《三国志》又有魏、蜀、吴不同的叙事位置。二五年和二二〇年的仪式材料都应被认真对待，因为仪式组织名分和官署；但不能让它们单独证明民意、军权或地方服从。兵数、屠杀规模、流民数量和具体行军路线尤其须保留材料的限度。
\end{sourcenote}
